\documentclass[9pt,t]{beamer}
% =====================================================================
% finance-beamer.tex — shared Beamer style for the LLM-in-Finance course
% =====================================================================
% Reusable slide style. Each deck \input's this immediately after
% \documentclass{beamer}. It provides:
%   * the Madrid theme + book colour palette (matches the printed book)
%   * two book-style framing blocks:
%       \begin{bigpicture} ... \end{bigpicture}   ("the bigger picture")
%       \begin{underhood}   ... \end{underhood}    ("under the hood")
%     These mirror the book's `context` and `deepdive` tcolorboxes so the
%     same intuition-vs-mechanism scaffold the reader meets in the book
%     reappears on the slides.
%   * a Python listings style for code frames
%   * appendix conventions: \appendixstart drops a divider and switches
%     to non-numbered backup frames so "extra / less central" material
%     lives after the main talk without inflating the slide count.
%
% Usage:
%   \documentclass{beamer}
%   % =====================================================================
% finance-beamer.tex — shared Beamer style for the LLM-in-Finance course
% =====================================================================
% Reusable slide style. Each deck \input's this immediately after
% \documentclass{beamer}. It provides:
%   * the Madrid theme + book colour palette (matches the printed book)
%   * two book-style framing blocks:
%       \begin{bigpicture} ... \end{bigpicture}   ("the bigger picture")
%       \begin{underhood}   ... \end{underhood}    ("under the hood")
%     These mirror the book's `context` and `deepdive` tcolorboxes so the
%     same intuition-vs-mechanism scaffold the reader meets in the book
%     reappears on the slides.
%   * a Python listings style for code frames
%   * appendix conventions: \appendixstart drops a divider and switches
%     to non-numbered backup frames so "extra / less central" material
%     lives after the main talk without inflating the slide count.
%
% Usage:
%   \documentclass{beamer}
%   % =====================================================================
% finance-beamer.tex — shared Beamer style for the LLM-in-Finance course
% =====================================================================
% Reusable slide style. Each deck \input's this immediately after
% \documentclass{beamer}. It provides:
%   * the Madrid theme + book colour palette (matches the printed book)
%   * two book-style framing blocks:
%       \begin{bigpicture} ... \end{bigpicture}   ("the bigger picture")
%       \begin{underhood}   ... \end{underhood}    ("under the hood")
%     These mirror the book's `context` and `deepdive` tcolorboxes so the
%     same intuition-vs-mechanism scaffold the reader meets in the book
%     reappears on the slides.
%   * a Python listings style for code frames
%   * appendix conventions: \appendixstart drops a divider and switches
%     to non-numbered backup frames so "extra / less central" material
%     lives after the main talk without inflating the slide count.
%
% Usage:
%   \documentclass{beamer}
%   \input{../_style/finance-beamer.tex}
%   \title{...}\subtitle{...}\author{...}\institute{...}\date{\today}
%   \begin{document} ... \end{document}
% =====================================================================

\usetheme{Madrid}
\usecolortheme{whale}
\setbeamertemplate{navigation symbols}{}
\setbeamertemplate{caption}[numbered]
\setbeamercovered{transparent}

% --- Density controls: keep dense, equation-heavy frames on one slide ---
% Slightly smaller body text + tighter lists/blocks. Frame titles and the
% Madrid chrome keep their normal size, so the deck still reads as Madrid,
% just with more vertical room for content.
\setbeamerfont{normal text}{size=\small}
\setbeamertemplate{itemize/enumerate body begin}{\small}
\setbeamertemplate{itemize/enumerate subbody begin}{\footnotesize}
% Tighter block spacing.
\setbeamertemplate{block begin}{%
  \par\vskip2pt%
  \begin{beamercolorbox}[colsep*=.75ex,leftskip=1ex,rightskip=1ex]{block title}%
    \usebeamerfont*{block title}\insertblocktitle%
  \end{beamercolorbox}%
  {\parskip0pt\vskip-1pt}%
  \begin{beamercolorbox}[colsep*=.75ex,vmode,leftskip=1ex,rightskip=1ex]{block body}%
  \vskip1pt}
\setbeamertemplate{block end}{\end{beamercolorbox}\vskip2pt}

\usepackage{amsmath, amssymb}
\usepackage{booktabs}
\usepackage{xcolor}
\usepackage{listings}
\usepackage[skins, breakable]{tcolorbox}

% --- Book colour palette (kept in sync with book/preamble.tex) ---
\definecolor{linkblue}{RGB}{14, 75, 140}
\definecolor{deepdivebg}{RGB}{235, 244, 255}
\definecolor{narrativebg}{RGB}{255, 252, 240}
\definecolor{narrativerule}{RGB}{180, 130, 40}
\definecolor{steelblue}{RGB}{70, 130, 180}
\definecolor{forestgreen}{RGB}{34, 139, 34}
\definecolor{crimson}{RGB}{220, 20, 60}
\definecolor{darkorange}{RGB}{255, 140, 0}
\definecolor{codebg}{RGB}{248, 248, 248}
\definecolor{coderule}{RGB}{210, 210, 210}
\definecolor{keywordblue}{RGB}{0, 100, 180}
\definecolor{commentgray}{RGB}{120, 120, 120}
\definecolor{stringgreen}{RGB}{30, 130, 30}

% Tie the Madrid structural colour to the book's link blue.
\setbeamercolor{structure}{fg=linkblue}
\setbeamercolor{title}{fg=white}
\setbeamercolor{frametitle}{fg=white}

% --- "the bigger picture" : narrative / intuition framing -----------
\newtcolorbox{bigpicture}{%
  enhanced, breakable, boxsep=2pt,
  colback  = narrativebg,
  colframe = narrativerule,
  boxrule  = 0pt, toprule = 1.4pt, bottomrule = 0.4pt,
  leftrule = 0pt, rightrule = 0pt, arc = 0pt,
  left = 6pt, right = 6pt, top = 4pt, bottom = 5pt,
  fonttitle = \footnotesize\scshape\color{narrativerule},
  title = {the bigger picture}, attach title to upper = {\par\smallskip},
}

% --- "under the hood" : the mechanism / technical detail ------------
\newtcolorbox{underhood}{%
  enhanced, breakable, boxsep=2pt,
  colback  = deepdivebg,
  colframe = linkblue,
  boxrule  = 0pt, toprule = 1.4pt, bottomrule = 0.4pt,
  leftrule = 0pt, rightrule = 0pt, arc = 0pt,
  left = 6pt, right = 6pt, top = 4pt, bottom = 5pt,
  fonttitle = \footnotesize\scshape\color{linkblue},
  title = {under the hood}, attach title to upper = {\par\smallskip},
}

% --- Python code style ---------------------------------------------
\lstset{
  basicstyle    = \ttfamily\footnotesize,
  keywordstyle  = \color{keywordblue}\bfseries,
  commentstyle  = \color{commentgray}\itshape,
  stringstyle   = \color{stringgreen},
  showstringspaces = false,
  language      = Python,
  breaklines    = true,
  backgroundcolor = \color{codebg},
  frame         = single,
  rulecolor     = \color{coderule},
  numbers       = none,
  aboveskip     = 4pt, belowskip = 4pt,
}

% --- Appendix conventions ------------------------------------------
% \appendixstart{Title} : begins the "extra material" appendix. Frames
% after it are not counted toward the main deck's frame total, keeping
% the core talk lean while preserving the deeper / optional content.
\newcommand{\appendixstart}[1]{%
  \appendix
  \setbeamertemplate{frametitle continuation}{}%
  \begin{frame}[noframenumbering]
    \begin{center}
      {\Large\scshape\color{linkblue} Appendix}\\[4pt]
      {\large #1}\\[10pt]
      {\footnotesize Extra / optional material — beyond the core lecture.}
    \end{center}
  \end{frame}
}
% Use \begin{frame}[noframenumbering]{...} for each appendix frame.

%   \title{...}\subtitle{...}\author{...}\institute{...}\date{\today}
%   \begin{document} ... \end{document}
% =====================================================================

\usetheme{Madrid}
\usecolortheme{whale}
\setbeamertemplate{navigation symbols}{}
\setbeamertemplate{caption}[numbered]
\setbeamercovered{transparent}

% --- Density controls: keep dense, equation-heavy frames on one slide ---
% Slightly smaller body text + tighter lists/blocks. Frame titles and the
% Madrid chrome keep their normal size, so the deck still reads as Madrid,
% just with more vertical room for content.
\setbeamerfont{normal text}{size=\small}
\setbeamertemplate{itemize/enumerate body begin}{\small}
\setbeamertemplate{itemize/enumerate subbody begin}{\footnotesize}
% Tighter block spacing.
\setbeamertemplate{block begin}{%
  \par\vskip2pt%
  \begin{beamercolorbox}[colsep*=.75ex,leftskip=1ex,rightskip=1ex]{block title}%
    \usebeamerfont*{block title}\insertblocktitle%
  \end{beamercolorbox}%
  {\parskip0pt\vskip-1pt}%
  \begin{beamercolorbox}[colsep*=.75ex,vmode,leftskip=1ex,rightskip=1ex]{block body}%
  \vskip1pt}
\setbeamertemplate{block end}{\end{beamercolorbox}\vskip2pt}

\usepackage{amsmath, amssymb}
\usepackage{booktabs}
\usepackage{xcolor}
\usepackage{listings}
\usepackage[skins, breakable]{tcolorbox}

% --- Book colour palette (kept in sync with book/preamble.tex) ---
\definecolor{linkblue}{RGB}{14, 75, 140}
\definecolor{deepdivebg}{RGB}{235, 244, 255}
\definecolor{narrativebg}{RGB}{255, 252, 240}
\definecolor{narrativerule}{RGB}{180, 130, 40}
\definecolor{steelblue}{RGB}{70, 130, 180}
\definecolor{forestgreen}{RGB}{34, 139, 34}
\definecolor{crimson}{RGB}{220, 20, 60}
\definecolor{darkorange}{RGB}{255, 140, 0}
\definecolor{codebg}{RGB}{248, 248, 248}
\definecolor{coderule}{RGB}{210, 210, 210}
\definecolor{keywordblue}{RGB}{0, 100, 180}
\definecolor{commentgray}{RGB}{120, 120, 120}
\definecolor{stringgreen}{RGB}{30, 130, 30}

% Tie the Madrid structural colour to the book's link blue.
\setbeamercolor{structure}{fg=linkblue}
\setbeamercolor{title}{fg=white}
\setbeamercolor{frametitle}{fg=white}

% --- "the bigger picture" : narrative / intuition framing -----------
\newtcolorbox{bigpicture}{%
  enhanced, breakable, boxsep=2pt,
  colback  = narrativebg,
  colframe = narrativerule,
  boxrule  = 0pt, toprule = 1.4pt, bottomrule = 0.4pt,
  leftrule = 0pt, rightrule = 0pt, arc = 0pt,
  left = 6pt, right = 6pt, top = 4pt, bottom = 5pt,
  fonttitle = \footnotesize\scshape\color{narrativerule},
  title = {the bigger picture}, attach title to upper = {\par\smallskip},
}

% --- "under the hood" : the mechanism / technical detail ------------
\newtcolorbox{underhood}{%
  enhanced, breakable, boxsep=2pt,
  colback  = deepdivebg,
  colframe = linkblue,
  boxrule  = 0pt, toprule = 1.4pt, bottomrule = 0.4pt,
  leftrule = 0pt, rightrule = 0pt, arc = 0pt,
  left = 6pt, right = 6pt, top = 4pt, bottom = 5pt,
  fonttitle = \footnotesize\scshape\color{linkblue},
  title = {under the hood}, attach title to upper = {\par\smallskip},
}

% --- Python code style ---------------------------------------------
\lstset{
  basicstyle    = \ttfamily\footnotesize,
  keywordstyle  = \color{keywordblue}\bfseries,
  commentstyle  = \color{commentgray}\itshape,
  stringstyle   = \color{stringgreen},
  showstringspaces = false,
  language      = Python,
  breaklines    = true,
  backgroundcolor = \color{codebg},
  frame         = single,
  rulecolor     = \color{coderule},
  numbers       = none,
  aboveskip     = 4pt, belowskip = 4pt,
}

% --- Appendix conventions ------------------------------------------
% \appendixstart{Title} : begins the "extra material" appendix. Frames
% after it are not counted toward the main deck's frame total, keeping
% the core talk lean while preserving the deeper / optional content.
\newcommand{\appendixstart}[1]{%
  \appendix
  \setbeamertemplate{frametitle continuation}{}%
  \begin{frame}[noframenumbering]
    \begin{center}
      {\Large\scshape\color{linkblue} Appendix}\\[4pt]
      {\large #1}\\[10pt]
      {\footnotesize Extra / optional material — beyond the core lecture.}
    \end{center}
  \end{frame}
}
% Use \begin{frame}[noframenumbering]{...} for each appendix frame.

%   \title{...}\subtitle{...}\author{...}\institute{...}\date{\today}
%   \begin{document} ... \end{document}
% =====================================================================

\usetheme{Madrid}
\usecolortheme{whale}
\setbeamertemplate{navigation symbols}{}
\setbeamertemplate{caption}[numbered]
\setbeamercovered{transparent}

% --- Density controls: keep dense, equation-heavy frames on one slide ---
% Slightly smaller body text + tighter lists/blocks. Frame titles and the
% Madrid chrome keep their normal size, so the deck still reads as Madrid,
% just with more vertical room for content.
\setbeamerfont{normal text}{size=\small}
\setbeamertemplate{itemize/enumerate body begin}{\small}
\setbeamertemplate{itemize/enumerate subbody begin}{\footnotesize}
% Tighter block spacing.
\setbeamertemplate{block begin}{%
  \par\vskip2pt%
  \begin{beamercolorbox}[colsep*=.75ex,leftskip=1ex,rightskip=1ex]{block title}%
    \usebeamerfont*{block title}\insertblocktitle%
  \end{beamercolorbox}%
  {\parskip0pt\vskip-1pt}%
  \begin{beamercolorbox}[colsep*=.75ex,vmode,leftskip=1ex,rightskip=1ex]{block body}%
  \vskip1pt}
\setbeamertemplate{block end}{\end{beamercolorbox}\vskip2pt}

\usepackage{amsmath, amssymb}
\usepackage{booktabs}
\usepackage{xcolor}
\usepackage{listings}
\usepackage[skins, breakable]{tcolorbox}

% --- Book colour palette (kept in sync with book/preamble.tex) ---
\definecolor{linkblue}{RGB}{14, 75, 140}
\definecolor{deepdivebg}{RGB}{235, 244, 255}
\definecolor{narrativebg}{RGB}{255, 252, 240}
\definecolor{narrativerule}{RGB}{180, 130, 40}
\definecolor{steelblue}{RGB}{70, 130, 180}
\definecolor{forestgreen}{RGB}{34, 139, 34}
\definecolor{crimson}{RGB}{220, 20, 60}
\definecolor{darkorange}{RGB}{255, 140, 0}
\definecolor{codebg}{RGB}{248, 248, 248}
\definecolor{coderule}{RGB}{210, 210, 210}
\definecolor{keywordblue}{RGB}{0, 100, 180}
\definecolor{commentgray}{RGB}{120, 120, 120}
\definecolor{stringgreen}{RGB}{30, 130, 30}

% Tie the Madrid structural colour to the book's link blue.
\setbeamercolor{structure}{fg=linkblue}
\setbeamercolor{title}{fg=white}
\setbeamercolor{frametitle}{fg=white}

% --- "the bigger picture" : narrative / intuition framing -----------
\newtcolorbox{bigpicture}{%
  enhanced, breakable, boxsep=2pt,
  colback  = narrativebg,
  colframe = narrativerule,
  boxrule  = 0pt, toprule = 1.4pt, bottomrule = 0.4pt,
  leftrule = 0pt, rightrule = 0pt, arc = 0pt,
  left = 6pt, right = 6pt, top = 4pt, bottom = 5pt,
  fonttitle = \footnotesize\scshape\color{narrativerule},
  title = {the bigger picture}, attach title to upper = {\par\smallskip},
}

% --- "under the hood" : the mechanism / technical detail ------------
\newtcolorbox{underhood}{%
  enhanced, breakable, boxsep=2pt,
  colback  = deepdivebg,
  colframe = linkblue,
  boxrule  = 0pt, toprule = 1.4pt, bottomrule = 0.4pt,
  leftrule = 0pt, rightrule = 0pt, arc = 0pt,
  left = 6pt, right = 6pt, top = 4pt, bottom = 5pt,
  fonttitle = \footnotesize\scshape\color{linkblue},
  title = {under the hood}, attach title to upper = {\par\smallskip},
}

% --- Python code style ---------------------------------------------
\lstset{
  basicstyle    = \ttfamily\footnotesize,
  keywordstyle  = \color{keywordblue}\bfseries,
  commentstyle  = \color{commentgray}\itshape,
  stringstyle   = \color{stringgreen},
  showstringspaces = false,
  language      = Python,
  breaklines    = true,
  backgroundcolor = \color{codebg},
  frame         = single,
  rulecolor     = \color{coderule},
  numbers       = none,
  aboveskip     = 4pt, belowskip = 4pt,
}

% --- Appendix conventions ------------------------------------------
% \appendixstart{Title} : begins the "extra material" appendix. Frames
% after it are not counted toward the main deck's frame total, keeping
% the core talk lean while preserving the deeper / optional content.
\newcommand{\appendixstart}[1]{%
  \appendix
  \setbeamertemplate{frametitle continuation}{}%
  \begin{frame}[noframenumbering]
    \begin{center}
      {\Large\scshape\color{linkblue} Appendix}\\[4pt]
      {\large #1}\\[10pt]
      {\footnotesize Extra / optional material — beyond the core lecture.}
    \end{center}
  \end{frame}
}
% Use \begin{frame}[noframenumbering]{...} for each appendix frame.


\title{Practical Session 2: Transformers, APIs, and RAG}
\subtitle{Hands-On: Attention by Hand, Cost Estimation, and a Minimal RAG Pipeline}
\author{Juan F.\ Imbet}
\institute{Paris Dauphine -- PSL University}
\date{Large Language Models in Finance}

\begin{document}

\begin{frame}\titlepage\end{frame}

% ============================================================
% FRAME 1: Overview
% ============================================================
\begin{frame}
\frametitle{Session Overview}

\textbf{Duration:} 1 hour (50 min problems + 10 min debrief)

\medskip
\textbf{Three activities:}
\begin{enumerate}
  \item \textbf{Problem 1} (15 min): Trace scaled dot-product attention by hand on a
        3-token financial sequence. Compute $QK^\top$, scale, softmax, multiply $V$.
  \item \textbf{Problem 2} (15 min): Estimate API costs for a realistic compliance
        workload. Compare GPT-4o, Claude 3 Haiku, and self-hosted Llama 3.
  \item \textbf{Case Study} (20 min): Run a minimal RAG pipeline in Python using
        \texttt{sentence\_transformers}, \texttt{faiss}, and \texttt{anthropic}.
        Answer three diagnostic questions about retrieval behaviour.
\end{enumerate}

\medskip
\begin{block}{Discussion (10 min)}
When would you choose BERT vs.\ GPT vs.\ RAG for a financial NLP task?
Give a concrete scenario for each.
\end{block}
\end{frame}

% ============================================================
% FRAME 2: Recap — Attention Formula
% ============================================================
\begin{frame}
\frametitle{Recap: Scaled Dot-Product Attention}

You will need this formula for Problem 1.

\medskip
\textbf{Given:} query matrix $Q$, key matrix $K$, value matrix $V$, head dimension $d_k$.

\[
\text{Attention}(Q, K, V) = \text{softmax}\!\left(\frac{QK^\top}{\sqrt{d_k}}\right) V
\]

\medskip
\textbf{Steps:}
\begin{enumerate}
  \item Compute $S = QK^\top$ \hfill (raw score matrix, shape $n \times n$)
  \item Scale: $\tilde{S} = S / \sqrt{d_k}$
  \item Apply softmax row-wise to get attention weights $A$
  \item Output: $O = AV$
\end{enumerate}

\medskip
\textbf{Softmax reminder:}
\[
\text{softmax}(\mathbf{x})_i = \frac{e^{x_i}}{\sum_j e^{x_j}}
\]

Numerical stability tip: subtract $\max(\mathbf{x})$ before exponentiating.
\end{frame}

% ============================================================
% FRAME 3: Recap — BPE and Cost Model
% ============================================================
\begin{frame}
\frametitle{Recap: BPE Tokenisation and Cost Formula}

You will need this for Problem 2.

\medskip
\textbf{Token conversion rule of thumb:}
\[
\text{tokens} \approx 1.3 \times \text{words}
\]

\medskip
\textbf{API cost model:}
\[
\text{cost} = N_{\text{in}} \cdot c_{\text{in}} + N_{\text{out}} \cdot c_{\text{out}}
\]

where $N_{\text{in}}$, $N_{\text{out}}$ are token counts and $c_{\text{in}}$, $c_{\text{out}}$ are
per-token prices in USD per million tokens (MTok).

\medskip
\begin{tabular}{lll}
\toprule
\textbf{Model} & \textbf{Input (\$/MTok)} & \textbf{Output (\$/MTok)} \\
\midrule
GPT-4o           & \$5.00 & \$15.00 \\
Claude 3 Haiku   & \$0.25 & \$1.25 \\
Llama 3 (A100)   & \$0 (compute only) & \$0 \\
\bottomrule
\end{tabular}

\medskip
\textbf{Self-hosted cost:} an A100 GPU costs approximately \$3/hr on cloud spot markets.
\end{frame}

% ============================================================
% FRAME 4: Problem 1 — Setup
% ============================================================
\begin{frame}
\frametitle{Problem 1: Attention by Hand — Setup}

\textbf{Sequence:} 3 tokens from a financial sentence: \texttt{"revenue"}, \texttt{"increased"}, \texttt{"significantly"}.

\medskip
\textbf{Head dimension:} $d_k = 2$.

\medskip
\textbf{Given matrices} (small integers for tractability):

\[
Q = \begin{pmatrix} 1 & 0 \\ 2 & 1 \\ 0 & 1 \end{pmatrix},
\quad
K = \begin{pmatrix} 1 & 1 \\ 0 & 2 \\ 1 & 0 \end{pmatrix},
\quad
V = \begin{pmatrix} 1 & 0 \\ 0 & 1 \\ 1 & 1 \end{pmatrix}
\]

Rows of $Q$, $K$, $V$ correspond to tokens \texttt{revenue} (row 1),
\texttt{increased} (row 2), \texttt{significantly} (row 3).

\medskip
\begin{block}{Interpretation}
$d_k = 2$ is unrealistically small; in GPT-3 $d_k = 64$; in GPT-4 $d_k \approx 128$.
The algebra is identical — we just keep arithmetic manageable.
\end{block}
\end{frame}

% ============================================================
% FRAME 5: Problem 1 — Tasks
% ============================================================
\begin{frame}
\frametitle{Problem 1: Tasks}

Complete the following steps. Show your working.

\medskip
\textbf{Task A.} Compute $S = QK^\top$ (a $3 \times 3$ matrix).\\
Recall: entry $(i,j) = \mathbf{q}_i^\top \mathbf{k}_j$.

\medskip
\textbf{Task B.} Scale by $\sqrt{d_k} = \sqrt{2} \approx 1.414$.\\
Report $\tilde{S} = S / \sqrt{2}$ to two decimal places.

\medskip
\textbf{Task C.} Apply softmax row-wise to $\tilde{S}$ to obtain the attention weight
matrix $A$ (each row sums to 1). Round to three decimal places.

\medskip
\textbf{Task D.} Compute the output $O = AV$.\\
Each row of $O$ is the attention-weighted sum of value vectors.

\medskip
\textbf{Discussion questions:}
\begin{itemize}
  \item Which value vector does token \texttt{revenue} attend to most heavily?
  \item What would happen to the attention weights if $d_k = 64$ instead of 2,
        with the same raw dot products? Would they be sharper or flatter?
\end{itemize}
\end{frame}

% ============================================================
% FRAME 6: Problem 2 — API Cost Estimation
% ============================================================
\begin{frame}
\frametitle{Problem 2: API Cost Estimation}

\textbf{Scenario:} A compliance team at a large European bank needs to process
$N = 5{,}000$ annual report filings (10-K equivalents). Each filing has
approximately 80,000 words. The pipeline uses:
\begin{itemize}
  \item A system prompt of 1,000 tokens (analytical instructions, fixed)
  \item Full filing text as input (converted to tokens)
  \item 600 output tokens per filing (structured compliance summary)
\end{itemize}

\medskip
\textbf{Task A.} Compute the total cost of processing all 5,000 filings using
\textbf{GPT-4o} (\$5/MTok input, \$15/MTok output).

\medskip
\textbf{Task B.} Compute the same cost using \textbf{Claude 3 Haiku}
(\$0.25/MTok input, \$1.25/MTok output).

\medskip
\textbf{Task C.} The team considers self-hosting Llama 3 on an A100 GPU rented
at \$3/hour. The pipeline can process 200 filings per hour. What is the total
compute cost? Under what conditions would self-hosting be preferred?

\medskip
\textbf{Bonus:} If the system prompt is cacheable and Anthropic offers 90\% discount
on cached input tokens, recompute the Claude 3 Haiku cost.
\end{frame}

% ============================================================
% FRAME 7: Case Study — Introduction
% ============================================================
\begin{frame}
\frametitle{Case Study: A Minimal RAG Pipeline for SEC Filings}

\textbf{Goal:} build a working RAG pipeline that retrieves relevant passages from
a small corpus of SEC filing excerpts and uses Claude to answer a question.

\medskip
\textbf{Stack:}
\begin{itemize}
  \item \texttt{sentence\_transformers} — bi-encoder embeddings (SBERT)
  \item \texttt{faiss} — approximate nearest-neighbour index
  \item \texttt{anthropic} — generation with the Claude API
\end{itemize}

\medskip
\textbf{Corpus:} 5 SEC filing passages (pre-loaded in the code).

\medskip
\textbf{What you will do:}
\begin{enumerate}
  \item Run the provided Python code
  \item Observe which passage is retrieved for a given query
  \item Answer three diagnostic questions about retrieval behaviour
\end{enumerate}

\medskip
\begin{alertblock}{Prerequisites}
You need an \texttt{ANTHROPIC\_API\_KEY} environment variable set.\\
Install: \texttt{pip install sentence-transformers faiss-cpu anthropic}
\end{alertblock}
\end{frame}

% ============================================================
% FRAME 8: Case Study — Code
% ============================================================
\begin{frame}[fragile]
\frametitle{Case Study: The Code --- Corpus and Index}

{\footnotesize
\begin{lstlisting}
from sentence_transformers import SentenceTransformer
import faiss, numpy as np, anthropic

PASSAGES = [
    "Revenue for FY2023 was $4.2 billion, up 12% year-over-year.",
    "Operating expenses increased to $1.8 billion due to R&D investment.",
    "Net income declined 5% to $820 million, impacted by higher interest expense.",
    "The company holds $2.1 billion in cash and equivalents.",
    "Guidance for FY2024 projects revenue growth of 8-10%.",
]

model = SentenceTransformer("all-MiniLM-L6-v2")
embs = model.encode(PASSAGES, normalize_embeddings=True).astype("float32")
index = faiss.IndexFlatIP(embs.shape[1])
index.add(embs)
\end{lstlisting}
}
\end{frame}

% ============================================================
% FRAME 8 (cont.): Case Study — Code (query function)
% ============================================================
\begin{frame}[fragile]
\frametitle{Case Study: The Code --- Retrieval and Generation}

{\footnotesize
\begin{lstlisting}
def rag_query(query: str, k: int = 2) -> str:
    q_emb = model.encode([query], normalize_embeddings=True).astype("float32")
    _, ids = index.search(q_emb, k)
    context = "\n".join(f"[{i+1}] {PASSAGES[ids[0][i]]}" for i in range(k))
    client = anthropic.Anthropic()
    msg = client.messages.create(
        model="claude-haiku-4-5",
        max_tokens=256,
        system="Answer using only the provided passages. Cite passage numbers as [N].",
        messages=[{"role": "user",
                   "content": f"{context}\n\nQuestion: {query}"}],
        temperature=0.0,
    )
    return msg.content[0].text

print(rag_query("What was the net income?"))
\end{lstlisting}
}
\end{frame}

% ============================================================
% FRAME 9: Case Study — Questions
% ============================================================
\begin{frame}
\frametitle{Case Study: Diagnostic Questions}

Run the code for the query \texttt{"What was the net income?"} and then answer:

\medskip
\textbf{Q1.} Which passage(s) were retrieved (i.e., which indices $i$ did FAISS return)?
Does the answer from Claude correctly cite those passage numbers? Why or why not?

\medskip
\textbf{Q2.} Change $k = 2$ to $k = 1$. Does the answer change? Change it to $k = 5$
(all passages). How does including more context affect the quality and length of the
answer?
\end{frame}

% ============================================================
% FRAME 9 (cont.): Case Study — Questions
% ============================================================
\begin{frame}
\frametitle{Case Study: Diagnostic Questions (cont.)}

\textbf{Q3.} Change \texttt{temperature=0.0} to \texttt{temperature=0.7}. Run the
same query three times. Do you observe variation in the output? Is that variation
desirable for this task? What does this tell you about the appropriate temperature
for structured financial extraction?

\medskip
\textbf{Q4 (Hallucination).} Modify the code to skip retrieval and send only the raw
question to Claude with no context. Does the model still produce an answer?
\begin{itemize}
  \item What type of hallucination is this (intrinsic or extrinsic)?
  \item Restore retrieval. Which mitigation technique from the lecture does RAG implement?
  \item Name one additional mitigation you would apply for a production compliance pipeline.
\end{itemize}
\end{frame}

% ============================================================
% FRAME 10: Discussion
% ============================================================
\begin{frame}
\frametitle{Discussion: When to Choose BERT vs. GPT vs. RAG}

For each scenario below, identify which approach you would use and why.

\medskip
\begin{enumerate}
  \item A compliance officer needs to classify 50,000 news headlines per day as
        \emph{material} or \emph{non-material} for insider-trading monitoring purposes.
        Latency target: under 50 ms per headline. Training data: 10,000 labelled examples.

  \item A portfolio manager wants a system that can answer questions about any 10-K
        filing in the portfolio in real time, using the most recent version of each filing.
        The filings were not part of any model's training data.

  \item An investment bank needs to draft the ``Risk Factors'' section of a new bond
        prospectus given a structured data dump of the issuer's financial statements and
        a comparable set of five precedent transactions.
\end{enumerate}

\medskip
\begin{block}{Frame your answer around:}
Latency, available labelled data, generation vs.\ classification, need for up-to-date
information, cost, and auditability.
\end{block}
\end{frame}

% ============================================================
% FRAME 11: Solution — Problem 1
% ============================================================
\begin{frame}
\frametitle{Solution: Problem 1 — Attention by Hand}

\textbf{Task A:} $S = QK^\top$

\[
S = \begin{pmatrix}1&0\\2&1\\0&1\end{pmatrix}
    \begin{pmatrix}1&0&1\\1&2&0\end{pmatrix}
  = \begin{pmatrix}1&0&1\\3&2&2\\1&2&0\end{pmatrix}
\]

\textbf{Task B:} $\tilde{S} = S/\sqrt{2}$

\[
\tilde{S} \approx \begin{pmatrix}0.71&0.00&0.71\\2.12&1.41&1.41\\0.71&1.41&0.00\end{pmatrix}
\]
\end{frame}

% ============================================================
% FRAME 11 (cont.): Solution — Problem 1
% ============================================================
\begin{frame}
\frametitle{Solution: Problem 1 --- Attention by Hand (cont.)}

\textbf{Task C:} Softmax row-wise (row 1 example: $e^{0.71}+e^{0}+e^{0.71} = 2.034+1+2.034$)

\[
A \approx \begin{pmatrix}0.400&0.200&0.400\\0.561&0.220&0.220\\0.260&0.529&0.211\end{pmatrix}
\]

\textbf{Task D:} $O = AV$

\[
O \approx \begin{pmatrix}0.800&0.800\\0.781&0.439\\0.740&0.740\end{pmatrix}
\]

\textbf{Discussion:} Token \texttt{revenue} splits attention evenly between itself and
\texttt{significantly}. Larger $d_k$ makes softmax \emph{sharper} (larger pre-softmax magnitudes).
\end{frame}

% ============================================================
% FRAME 12: Solution — Problem 2
% ============================================================
\begin{frame}
\frametitle{Solution: Problem 2 — API Cost Estimation}

\textbf{Token counts per filing:}
\begin{itemize}
  \item Input: $80{,}000 \times 1.3 = 104{,}000$ tokens from filing text
  \item $+ 1{,}000$ system prompt $= 105{,}000$ input tokens per filing
  \item Output: $600$ tokens per filing
\end{itemize}

\medskip
\textbf{Task A — GPT-4o:}
\[
\text{cost} = 5{,}000 \times \left(\frac{105{,}000}{10^6} \times 5.00
+ \frac{600}{10^6} \times 15.00\right)
= 5{,}000 \times (0.525 + 0.009) = \$2{,}670
\]

\textbf{Task B — Claude 3 Haiku:}
\[
\text{cost} = 5{,}000 \times \left(\frac{105{,}000}{10^6} \times 0.25
+ \frac{600}{10^6} \times 1.25\right)
= 5{,}000 \times (0.02625 + 0.00075) = \$135
\]
\end{frame}

% ============================================================
% FRAME 12 (cont.): Solution — Problem 2
% ============================================================
\begin{frame}
\frametitle{Solution: Problem 2 --- API Cost Estimation (cont.)}

\textbf{Task C — Self-hosted Llama 3 (A100 at \$3/hr, 200 filings/hr):}
\[
\text{time} = 5{,}000/200 = 25 \text{ hr},\quad
\text{cost} = 25 \times 3 = \$75
\]

\textbf{Bonus — Prompt caching (90\% off system prompt for Claude Haiku):}
\[
\text{cost} \approx 5{,}000 \times \left(\frac{104{,}000}{10^6}{\cdot}0.25 + \frac{1{,}000}{10^6}{\cdot}0.025 + \frac{600}{10^6}{\cdot}1.25\right) \approx \$133
\]

\textbf{When to self-host:} data privacy requirements; $>$\$3k/yr equivalent spend; latency not critical.
\end{frame}

% ============================================================
% FRAME 13: Wrap-Up
% ============================================================
\begin{frame}
\frametitle{Wrap-Up: Key Takeaways from This Session}

\textbf{From Problem 1:}
\begin{itemize}
  \item Attention is a weighted average of value vectors; the weights come from
        (softmaxed) scaled dot products between queries and keys
  \item Larger $d_k$ sharpens the softmax --- which is why we normalise by $\sqrt{d_k}$
  \item The mechanism is fully differentiable and can be traced algebraically
\end{itemize}

\medskip
\textbf{From Problem 2:}
\begin{itemize}
  \item Cost differences between API providers can be 20--35$\times$ for the same workload
  \item Self-hosting is competitive when the workload is large and latency is flexible
  \item Prompt caching for fixed system prompts provides significant savings
\end{itemize}
\end{frame}

% ============================================================
% FRAME 13 (cont.): Wrap-Up
% ============================================================
\begin{frame}
\frametitle{Wrap-Up: Key Takeaways (cont.)}

\textbf{From the Case Study:}
\begin{itemize}
  \item RAG reliably retrieves semantically relevant passages, even from small corpora
  \item $k$ and temperature are meaningful hyperparameters with measurable output effects
  \item For structured financial extraction: $\tau = 0$, small $k$, grounded system prompt
\end{itemize}

\medskip
\begin{block}{Next practical}
Practical Session 3: Fine-tuning FinBERT on FinancialPhraseBank and comparing
zero-shot GPT-4o against a fine-tuned encoder model on the all-agree split.
\end{block}
\end{frame}

% ============================================================
\appendixstart{Stretch Problem}
% ============================================================

% --- Optional extension (relocated from the Case Study frame) ---
\begin{frame}[noframenumbering]{Stretch Problem: Targeting a Single Passage}

\begin{block}{Optional extension}
Add a sixth passage: \texttt{"The company's debt-to-equity ratio is 1.8 as of year-end."}
Formulate a query that retrieves \emph{only} this passage. What makes a query specific
enough to target a single passage in a dense retrieval system?
\end{block}

\medskip
\textbf{Things to consider:}
\begin{itemize}
  \item How distinctive is the vocabulary of this passage (``debt-to-equity'',
        ``ratio'') relative to the other five?
  \item Does adding numeric or lexical anchors (e.g.\ ``1.8'', ``year-end'')
        sharpen the dense similarity score for this passage?
  \item Where would BM25 (sparse) help target an exact term that SBERT blurs?
\end{itemize}
\end{frame}

\end{document}
